\chapter{Panoramica dell'Architettura del Software}

\section{Analisi dei Componenti dell'Architettura a Livelli}
La piattaforma \textit{2Chance} si basa su un'architettura modulare a livelli che segue il tradizionale paradigma Model--View--Controller (MVC). Di conseguenza, le responsabilità di sistema sono rigorosamente mappate: la presentazione dei dati è gestita dalla View (JavaServer Pages e JavaScript); il coordinamento degli input è gestito dal Controller; e il Model è suddiviso in Service (logica di business), Data Access Object (persistenza) e Bean (strutture dati). Questo rigoroso disaccoppiamento costituisce la base esistente su cui sono stati applicati i miglioramenti di dependability del progetto.

\subsection{Il Livello View (Presentazione e Interazione)}
\noindent\textbf{Componente:} pacchetto \texttt{webapp} (es. \texttt{index.jsp}, \texttt{login.jsp}).

\noindent Il livello View costituisce la logica di presentazione ed è responsabile del rendering dell'interfaccia utente (UI) e dell'acquisizione degli input dell'utente. È implementato utilizzando i linguaggi standard del web, tra cui HyperText Markup Language (HTML), Cascading Style Sheets (CSS) per lo stile e JavaScript (JS) per l'interattività lato client. Utilizzando JavaServer Pages (JSP), visualizza dinamicamente i dati forniti dal Controller, tipicamente passati tramite attributi di request o di sessione.\\

Sebbene la View avvii il flusso di dati trasmettendo i parametri dell'utente tramite richieste HTTP, le è architettonicamente proibito eseguire logica di business o accedere direttamente al livello di persistenza.

\subsection{Il Livello Controller (Coordinamento degli Input e Routing)}
\noindent\textbf{Componente:} pacchetto \texttt{controller} (es. \texttt{LoginServlet}, \texttt{RegistrazioneServlet}).

\noindent Il Controller funge da punto di ingresso centralizzato per l'elaborazione lato server. Implementato tramite Java Servlet, intercetta e gestisce le richieste HTTP in ingresso (nello specifico, \texttt{GET} e \texttt{POST}) trasmesse dal client. Agisce come direttore del traffico dell'applicazione, assicurando che le azioni dell'utente vengano instradate agli appropriati gestori della logica di business. Il Controller agisce strettamente come intermediario e non contiene logica di business.\\

Il suo flusso di lavoro operativo consiste in tre fasi distinte:
\begin{enumerate}
	\item \textbf{Parsing:} estrae i parametri di input dall'oggetto request e gestisce la sessione HTTP per preservare lo stato dell'utente attraverso interazioni stateless.
	\item \textbf{Delegation (Delega):} invoca i metodi appropriati all'interno del livello Service, passando solo gli argomenti necessari.
	\item \textbf{Routing (Instradamento):} al completamento dell'esecuzione del servizio, determina la strategia di risposta, inoltrando (forward) il contesto della richiesta a una vista JSP per il rendering o reindirizzando (redirect) l'utente verso una nuova risorsa per prevenire il reinvio del modulo.
\end{enumerate}

\subsection{Il Livello Service (Orchestrazione della Logica di Business)}
\noindent\textbf{Componente:} pacchetto \texttt{services} (es. \texttt{LoginService}, \texttt{AdminService}).

\noindent Il livello Service costituisce il nucleo operativo dell'applicazione, incapsulando la logica di business del sistema. Serve come astrazione che disaccoppia la gestione degli input (Controller) dall'archiviazione dei dati (DAO). Isolando queste responsabilità, l'architettura centralizza le regole di business e le rende indipendenti dall'interfaccia utente e dalla tecnologia del database.\\

Implementati come classi Java standard, i componenti Service orchestrano flussi di lavoro complessi e fanno rispettare gli invarianti di dominio. In particolare, un Service convalida gli input rispetto alle regole specifiche dell'applicazione (es. assicurando che una password soddisfi i criteri di complessità o verificando la disponibilità in magazzino) prima di intraprendere qualsiasi azione. Per preservare la purezza architetturale, al livello Service è proibito eseguire direttamente query SQL. Invece, coordina la persistenza invocando il livello Data Access per recuperare o persistere i dati, garantendo che l'elaborazione logica rimanga distinta dalla manipolazione fisica dei dati.

\subsection{Il Livello Data Access (Persistenza e DAO)}
\noindent\textbf{Componente:} pacchetto \texttt{model.dao} (es. \texttt{UtenteDAO}, \texttt{ProdottoDAO}).

\noindent Il livello Data Access implementa il pattern Data Access Object (DAO), che astrae e incapsula tutto l'accesso alla fonte dei dati. Questo design assicura che il resto dell'applicazione rimanga agnostico rispetto all'implementazione concreta del database, consentendo potenziali modifiche nel meccanismo di archiviazione senza avere un impatto sulla logica di livello superiore. Questo livello è l'unica autorità autorizzata a eseguire comandi SQL.\\

Utilizzando JDBC (Java Database Connectivity) e un connection pool di Tomcat per un'efficiente gestione delle risorse, i DAO gestiscono la comunicazione a basso livello con il database. I loro compiti principali includono:
\begin{enumerate}
	\item \textbf{Operazioni CRUD:} esecuzione di transazioni Create, Read, Update e Delete.
	\item \textbf{Object--Relational Mapping (ORM):} trasformazione manuale dei dati relazionali grezzi (es. righe di \texttt{ResultSet}) in oggetti Java ad alto livello (Model Bean) e viceversa. Questa traduzione consente al livello Service di manipolare oggetti Java anziché tipi di dati specifici del database.
\end{enumerate}

\subsection{I Model Bean (Rappresentazione dei Dati)}
\noindent\textbf{Componente:} pacchetto \texttt{model.beans} (es. \texttt{Utente}, \texttt{Prodotto}).

\noindent I Model Bean rappresentano le entità funzionali fondamentali del sistema. Strutturalmente, sono implementati come Plain Old Java Object (POJO). La loro responsabilità principale è l'incapsulamento dei dati, consentendo il trasporto dello stato attraverso i livelli View, Controller, Service e Data Access.\\

Ogni Bean rispecchia una specifica entità nel modello di dominio, corrispondendo tipicamente a una tabella nel database relazionale. I Bean mantengono lo stato attraverso attributi privati, ai quali si accede e che si modificano esclusivamente tramite metodi pubblici di accesso (getter) e mutazione (setter).

\section{Comportamento Dinamico: Il Flusso di Autenticazione}
Per illustrare l'interazione dinamica tra i componenti architetturali, esaminiamo il flusso di controllo di una procedura standard di autenticazione utente (login). Questa sequenza mostra come i dati si propagano dall'interfaccia utente al database e viceversa, rispettando la separazione delle responsabilità (separation of concerns) definita sopra:
\begin{enumerate}
	\item \textbf{Interazione e richiesta (livello View):} l'utente inserisce le credenziali in \texttt{login.jsp}. All'invio, il browser genera una richiesta HTTP \texttt{POST} contenente i parametri (email e password) e la trasmette al server.
	\item \textbf{Intercettazione e delega (livello Controller):} la \texttt{LoginServlet} intercetta la richiesta e ne analizza il corpo per estrarre i parametri di input. In conformità con il suo ruolo di coordinamento, non esegue alcuna verifica dei dati; al contrario, delega l'operazione invocando \texttt{login()} su \texttt{LoginService}.
	\item \textbf{Orchestrazione e validazione (livello Service):} \texttt{LoginService} riceve i parametri grezzi ed esegue una validazione preliminare (es. assicurandosi che i campi non siano vuoti). Se la validazione ha esito positivo, richiede il recupero dei dati chiamando \texttt{getUserByEmailPassword()} su \texttt{UtenteDAO}.
	\item \textbf{Esecuzione della query (livello Data Access):} \texttt{UtenteDAO} stabilisce una connessione al database, costruisce ed esegue una query \texttt{SELECT} sicura. Questo è l'unico punto nel flusso in cui l'applicazione comunica direttamente con l'archiviazione di persistenza.
	\item \textbf{Mappatura degli oggetti (livello Data Access):} dal \texttt{ResultSet} restituito, il DAO esegue il mapping oggetto-relazionale: istanzia un nuovo Bean \texttt{Utente}, popola i suoi campi con i valori del database e restituisce l'oggetto al livello Service.
	\item \textbf{Gestione dello stato e risposta (livello Controller):} il Service restituisce il Bean \texttt{Utente} popolato al Controller. Il Controller quindi (i) persiste lo stato memorizzando il Bean nella sessione HTTP, effettuando di fatto l'accesso dell'utente, e (ii) esegue la navigazione inviando una risposta di reindirizzamento che istruisce il client a caricare \texttt{dashboard.jsp}, dove viene visualizzato il nome dell'utente.
\end{enumerate}

\section{Miglioramenti di Dependability e Toolchain}
In seguito allo sviluppo iniziale, l'architettura è stata significativamente modernizzata durante il progetto di software dependability. Questa fase ha introdotto diverse tecnologie e metodologie avanzate per garantire la qualità del codice, la sicurezza e la manutenibilità.

\subsection{Testing e Analisi della Copertura (JUnit, Mockito, JaCoCo e Codecov)}
È stata implementata una metodologia di testing completa e bottom-up per convalidare i componenti del sistema. \textbf{JUnit 5} funge da framework principale per i test automatizzati, mentre \textbf{Mockito} è ampiamente utilizzato per creare mock delle dipendenze esterne, come le connessioni al database nei DAO e le richieste HTTP nei Controller, consentendo unit test isolati. Inoltre, \textbf{JaCoCo} è integrato per misurare la copertura del codice (es. line e branch coverage), e i risultati vengono aggregati e visualizzati tramite \textbf{Codecov} durante la pipeline di Continuous Integration (CI), garantendo che gli standard di test siano continuamente rispettati.

\subsection{Specifica e Verifica Formale (JML e OpenJML)}
Per elevare l'affidabilità delle classi critiche, sono stati adottati metodi formali. Il \textbf{Java Modeling Language (JML)} è utilizzato per annotare il codice Java con specifiche formali, definendo precondizioni, postcondizioni e invarianti di classe. \textbf{OpenJML} viene quindi impiegato come strumento di verifica statica per dimostrare matematicamente che l'implementazione Java soddisfi rigorosamente tali specifiche, prevenendo errori logici e fallimenti in casi limite (edge-case).

\subsection{Mutation Testing (PITest)}
Per valutare l'effettiva capacità di rilevamento dei guasti delle suite di test, è stato introdotto il mutation testing tramite \textbf{PITest}. Invece di fare affidamento esclusivamente sulla line coverage, PITest inietta automaticamente piccoli guasti sintattici (mutanti) nel codice dell'applicazione ed esegue la suite di test. Se i test falliscono (uccidendo il mutante), ciò verifica che i test sono sufficientemente robusti per individuare le regressioni.

\subsection{Containerizzazione e Orchestrazione (Docker e Docker Compose)}
Per eliminare il problema del ``funziona sulla mia macchina'' e semplificare i deployment, l'infrastruttura dell'applicazione è stata containerizzata utilizzando \textbf{Docker}. \textbf{Docker Compose} è utilizzato per orchestrare l'ambiente multi-container, definendo servizi interconnessi distinti per l'applicazione web, il database MySQL e le utilità di testing. Questa configurazione garantisce un ambiente di esecuzione coerente e riproducibile attraverso lo sviluppo, le pipeline CI/CD e la produzione.

\subsection{Sicurezza e Valutazione delle Vulnerabilità}
È stata stabilita una solida pipeline di sicurezza automatizzata all'interno delle GitHub Actions per monitorare continuamente il repository e i suoi artefatti:
\begin{itemize}
    \item \textbf{GitGuardian (Secret Detection):} Implementato per la scansione automatizzata dei secret, analizza i commit per rilevare e prevenire l'esposizione di password codificate, chiavi API o token sensibili.
    \item \textbf{Snyk (Software Composition Analysis):} Impiegato per la SCA, scansiona attivamente le dipendenze Maven per identificare e mitigare le vulnerabilità di sicurezza note (CVE) nelle librerie di terze parti.
    \item \textbf{SonarQube Cloud (Static Analysis):} Utilizzato per il Static Application Security Testing (SAST), esegue una profonda analisi statica sulla base di codice per rilevare code smell, debito tecnico e vulnerabilità di sicurezza, in particolare all'interno del livello Controller.
\end{itemize}